\chapter{练习解答}

\section*{第\ref{ch:model}章}
1．一个数学上存在的随机变量未必属于可观察材料。让控制者读取真实种子或包含情境的变量，会改变可用信息与策略类别，从而取消原有下界的条件。

2．两个历史最后都得到“候选A仍可行”。第一个历史已经排除B，第二个尚未检查B。下一操作可以分别为细化A与检查B；仅看最后片段无法作出这种区分。

\section*{第\ref{ch:identifiability}章}
1．三个类分别贡献$0$、$0.3/2=0.15$与$0.3/3=0.1$，总误差为$0.25$。

2．添加实验只会为等价关系增加等同性约束，所以每个新类都是旧类的子集。可以忽略新增实验继续旧策略，故最优风险不增；也可直接比较类内最大先验质量。

\section*{第\ref{ch:coverage}章}
1．第一实验权重$w$时，两类覆盖分别为$aw$与$b(1-w)$。最优点令二者相等，故$w=b/(a+b)$，值为$ab/(a+b)$。

2．错误输出必有一个目标不同的情境在似然上击败真实情境，因此情境比较足以保证目标正确。对许多具有相同目标的情境逐个作并合界可能多付出$\ln m$；任务专用读出与表示有机会改进它。

\section*{第\ref{ch:lower}章}
1．使用附录中的两个偏导符号，将$p$降至$1-\delta$、将$q$升至$\delta$都不增大二元相对熵。

2．停止透露第一次观察属于预定停止集合。补空操作保留实际停止位置，完整记录分布仍包含这项信息；若删除位置，只比较存在的输出，便可能遗漏它。

\section*{第\ref{ch:operational}章}
1．例如$M$随规模指数增长，$C_M$也增长。带有$C_M$的下界可能变弱，不能再视为统一常数因子的匹配。

2．以$\lceil\ln(m/\delta)/u\rceil$次调用运行公开分配，可以保证最坏错误不超过$\delta$；另行报告分配取得、控制和通信费用。未知最优值不妨碍可执行上界，但不能据此声称最优。

\section*{第\ref{ch:calibration}章}
1．$(1,0,1)$是两行之和，可辨别。两行均与$(1,1,1)$正交，但目标$(1,1,1)$与该向量内积为1，因此后者不可辨别。

2．给目标涉及的$k$个基向量等权，任一目标冲突差至少命中其中一个坐标，得到$1/k$。反方等权选择这些单坐标差，任何实验的平均命中至多$1/k$，由对偶得最优。

\section*{第\ref{ch:communication}章}
1．所有情境的接收分布均为该符号的点质量，所以所有区别量为零，任意分配覆盖率都为零。

2．公平目标已被A知道，A仅在目标为1时发送固定标记。B看到沉默后应以概率1判断0；如果把沉默当作没有发生任何条件化，它会错误地保留先验$1/2$。

\section*{第\ref{ch:representation}章}
1．总分数误差不超过$E=\sum_je_j$，风险至多$\min\{1,me^E(1-u)^n\}$。若要求不超过$\delta$，一个充分条件为$E\le-n\ln(1-u)-\ln(m/\delta)$，右侧需非负。

2．可逆变换不合并情境似然区别，但逆变换可能很昂贵。例如将一个容易直接读取的答案嵌入需要大量计算才能解码的可逆表示，信息不丢失而受限读出变难。

\section*{第\ref{ch:interpretation}章}
1．成员独立产生噪声，但都观察$Y\oplus B$，共享持续偏差$B$。给定$(Y,B)$的独立性成立；相对于目标$Y$的共同误差并未消除。

2．固定目标冲突对，检查表示是否改变似然比或可辨别类；实现编码与读出并计费；再在全部有限情境或预先指定的实例族上检查风险，而非只复用构造表示时见过的答案。

\section*{第\ref{ch:frontiers}章}
1．第\ref{thm:dimension}条定理的$d=2^r$增长，而$M,\kap$固定。一般主定理的$d$因子变松；补充的信息下界恢复了关于$r$的正确增长率。

2．至少给出合法输入、观察协议、策略允许信息、成本、量词、构造、完整类别下界或反例、可复算数据，以及与最相近结果在假设和结论上的逐项差别。
